\chapter{Kết Luận}

Đề tài đã xây dựng được một mô hình mô phỏng RTL cho hệ thống CPU RISC-V 8 lõi theo kiến trúc pipeline 5 tầng. Các module chính đã được thiết kế và tích hợp gồm \texttt{risc\_core}, các thanh ghi đệm pipeline, cơ chế hazard/forwarding, \texttt{alu}, \texttt{pc\_logic}, \texttt{reg\_file}, \texttt{instr\_mem}, \texttt{data\_mem}, \texttt{arbiter}, \texttt{bus\_ctrl} và khối tích hợp toàn hệ thống \texttt{top\_8core}. Trên cơ sở đó, hệ thống có thể mô phỏng nhiều lõi cùng thực thi chương trình, phát sinh yêu cầu truy cập bộ nhớ dùng chung và được điều phối qua interconnect trung tâm.

Các chức năng đã được kiểm chứng gồm thực thi nhóm lệnh RV32I subset cơ bản, xử lý stall/flush trong pipeline, đọc/ghi shared memory, phân xử bus Round-Robin, phát hiện hệ thống không deadlock trong test tích hợp, cũng như kiểm chứng riêng các primitive đồng bộ gồm LR/SC success, LR/SC fail và AMOADD trong \texttt{interconnect\_tb}. Kết quả mô phỏng cho thấy các testbench chính đều đạt pass theo tiêu chí đã đặt ra, đồng thời tạo được waveform để quan sát các tín hiệu quan trọng bằng GTKWave.

Với phạm vi đồ án Kiến trúc máy tính, hệ thống đã đạt mục tiêu mô phỏng một CPU RISC-V đa lõi có pipeline và bộ nhớ chia sẻ ở mức RTL. Tuy nhiên, đây chưa phải là một bộ xử lý thương mại hoàn chỉnh. Các phần chưa được kiểm chứng đầy đủ gồm benchmark hiệu năng chuẩn, full RISC-V compliance, exception/interrupt/CSR, cache coherence, hệ điều hành và runtime đa luồng thực tế. Do đó, các kết quả hiện tại nên được hiểu là bằng chứng chức năng và kiến trúc cho mô hình nghiên cứu, không phải số liệu hiệu năng cuối cùng của một CPU hoàn chỉnh.

Giá trị học thuật chính của dự án nằm ở việc kết nối nhiều kiến thức nền của môn học vào một hệ thống chạy được: pipeline 5 tầng, hazard handling, datapath/control path, arbitration, shared-memory multicore và atomic primitive. Qua đó, báo cáo không chỉ mô tả từng module riêng lẻ mà còn chỉ ra cách các module phối hợp để hình thành một luồng thực thi hoàn chỉnh từ lệnh RISC-V đến giao dịch bộ nhớ và kết quả mô phỏng.
